\documentclass[aspectratio=169]{beamer}
\usepackage[utf8]{inputenc}
\usepackage[T1]{fontenc}
\usepackage[english]{babel}
\usepackage{hyperref}
\usepackage{tikz}
\usepackage{booktabs}
\usetikzlibrary{arrows.meta,positioning,shapes.geometric}

\hypersetup{colorlinks=true, urlcolor=blue, linkcolor=blue}

\newcommand{\pkg}[1]{\texttt{#1}}

\title{1.1.2 Development Phase / Reflection Phase}
\subtitle{Sentiment analysis from news articles - Project: Artificial Intelligence
(DLBDSEAIS02)
}
\author{Alejandro Moral Aranda}
\date{\href{https://github.com/alejandromoralwork/sentiment-analysis}{GitHub Link}}

\begin{document}

\begin{frame}
    \titlepage
\end{frame}

\begin{frame}{Short overview}
The conception from phase 1 was turned into a working tool. The app fetches news by keyword, gets the full article text, runs sentiment analysis with three models, and saves the result for later review. I kept the implementation small and practical, because the goal was a usable student project with reasonable performance.
\end{frame}

\begin{frame}{Frameworks and tools}
\begin{itemize}
    \item \href{https://fastapi.tiangolo.com/}{FastAPI} for the backend and WebSocket messages.
    \item \href{https://newsapi.org/}{NewsAPI} for recent news articles.
    \item \href{https://www.crummy.com/software/BeautifulSoup/bs4/doc/}{BeautifulSoup} for HTML text extraction.
    \item \href{https://www.nltk.org/vader.html}{VADER}, \href{https://textblob.readthedocs.io/}{TextBlob}, and \href{https://huggingface.co/docs/transformers/index}{Transformers} for sentiment analysis.
    \item \href{https://pandas.pydata.org/}{pandas} for CSV output.
    \item \href{https://www.chartjs.org/}{Chart.js} and \href{https://tailwindcss.com/}{Tailwind CSS} for the frontend.
\end{itemize}
\end{frame}

\begin{frame}{Setup of the project}
The first step was getting the app structure ready and installing the packages. The backend runs in app.py, and src/main.py is a local runner for testing the pipeline offline. The code is organized in small pieces for readability.

Important files:
\begin{itemize}
    \item app.py - FastAPI server and WebSocket handler
    \item src/news_fetcher.py - NewsAPI integration and text extraction
    \item src/sentiment.py - preprocessing and sentiment models
    \item src/reporting.py - saves results to CSV and JSON
    \item templates/index.html and templates/script.js - web dashboard
\end{itemize}
\end{frame}

\begin{frame}{Implementation flow}
The main flow works like this:

\begin{enumerate}
    \item User enters a keyword.
    \item Fetch articles from NewsAPI.
    \item Download the article page and extract clean text.
    \item Run VADER, TextBlob, and transformer models on the text.
    \item Take the consensus vote and calculate confidence.
    \item Save the result as JSON and CSV reports.
    \item Show live progress updates in the browser at each step.
\end{enumerate}
\end{frame}

\begin{frame}{Implemented components}
The diagram from the conception phase is now real code. The app does this:
\begin{enumerate}
    \item User gives a keyword.
    \item NewsAPI returns article links.
    \item The article HTML is downloaded and cleaned.
    \item The text is preprocessed to remove noise.
    \item All three models analyze the same text.
    \item The app stores a final consensus label.
    \item The report is saved and shown in the dashboard.
\end{enumerate}
The websocket part is useful because the user can see what is going on while the analysis is running.
\end{frame}

\begin{frame}{Code and models}
The code includes three sentiment analysis models:

\begin{itemize}
    \item \textbf{VADER} (Valence Aware Dictionary and sEntiment Reasoner) - lexicon-based approach.
    \item \textbf{TextBlob} - simple polarity and subjectivity scoring.
    \item \textbf{RoBERTa sentiment model} - neural model for context-aware analysis.
\end{itemize}

Each model analyzes the full article text. The app then votes on a consensus label and stores confidence scores. Comments in the code explain the key steps: text preprocessing, model execution, and report generation.
\end{frame}

\begin{frame}{Data and evaluation}
The project uses real news articles from NewsAPI, so the text is messy like actual news and not like a clean demo dataset. That helps a lot, because sentiment on news is hard and headlines alone are usually not enough.

\vspace{2mm}
For checking the output, I compare the models against each other and look at the report files that get written to \pkg{data/reports/}. If later a small manually labeled set is added, it can be used to measure accuracy, precision, recall, and F1.
\end{frame}

\begin{frame}{NLP models used}
\begin{itemize}
    \item \textbf{VADER}: fast baseline, good for short text, but not always strong on long news.
    \item \textbf{TextBlob}: easy to use, simple rule-style result.
    \item \textbf{RoBERTa sentiment model}: slower, but better at context and more useful for full articles.
\end{itemize}

\vspace{2mm}
The app does not rely on only one model. It uses all three and then takes the majority label. That helps when one model is too confident on a strange sentence.
\end{frame}

\begin{frame}{Iterative optimization}
The biggest fixes during development were pretty basic but important:
\begin{itemize}
    \item use the full article text instead of just the headline
    \item clean HTML before analysis
    \item keep the three-model comparison
    \item store reports so the last run is not lost
    \item show the output in the browser right away
\end{itemize}

\vspace{2mm}
These changes made the tool feel much more useful, because news text has a lot of noise and one model alone was giving too many wrong calls.
\end{frame}

\begin{frame}{Packaging, documentation and delivery}
The project was prepared for non-technical delivery to the marketing department:
\begin{itemize}
    \item A CLI entrypoint (`cli.py`) was added so the tool can be run from a command prompt.
    \item Packaging notes and a helper build script (`build_exe.bat`) allow creation of a single-file Windows executable with `pyinstaller`.
    \item Documentation: `docs/UserGuide.md` (for marketing users) and `docs/DeveloperGuide.md` (for maintainers) were included.
    \item Observability: `logging` replaced `print()` calls and a `requirements.pinned.txt` file was added for reproducible installs.
\end{itemize}
These changes make the deliverable more suitable for deployment and everyday use by non-technical stakeholders.
\end{frame}

\begin{frame}{Result and reflection}
At the end of this phase the project became a working news sentiment tool. It can fetch articles, clean the text, analyze the sentiment, and keep the result in a report. The code is still small and simple, but it now follows the plan from the conception phase much better.

\vspace{2mm}
The main lesson is that news sentiment is not easy. Full text works better than headlines, and comparing more than one model gives a more honest result.
\end{frame}



\end{document}